% Printed 401
\par\clearpage
\noindent\null\hfill{\small\textbf{ZUR HANNOVERSCHEN TRIANGULATION.}}\hfill\makebox[0pt][r]{\fontsize{13}{16}\selectfont 401}\par
\vspace*{4\baselineskip}
\begin{center}
{\large III.}\par\medskip
{\LARGE\textls[150]{NACHLASS}}\par\smallskip
{\large ZUR HANNOVERSCHEN TRIANGULATION.}
\end{center}
\vspace{2\baselineskip}
\begin{center}
[1.]\par\smallskip
{\large [Plan und Anfang zum Werke über die trigonometrischen Messungen\\
in Hannover.]}\par\medskip
\textls[120]{Plan des Werkes.}\par\smallskip
Acht Abschnitte.
\end{center}

\noindent
\begin{tabular*}{\linewidth}{@{}r@{\enspace}l@{\extracolsep{\fill}}r@{\hspace{3.6em}}}
1) & Über die bei den Messungen angewandten Einrichtungen & \\
   & und Methoden im allgemeinen & \makebox[0pt][r]{ungefähr\enspace}40\makebox[0pt][l]{\enspace Seiten}\\
2) & Messung der Hauptdreieckswinkel & 60\\
3) & Kleine Reductionen der Winkelmessungen & 12\\
4) & Ausgleichung der Messungen & 24\\
5) & Erste Berechnungsart der Dreiecke, in der Darstellung & \\
   & auf dem Sphäroid & 36\\
6) & Zweite Berechnungsart, in der Darstellung in der Ebene & 60\\
7) & Nebenpunkte & 24\\
8) & Höhenmessungen & 20\\
\multicolumn{3}{@{}r@{}}{Zusammen etwa 276 Seiten}\\
\multicolumn{3}{@{}r@{}}{oder \(34\tfrac{1}{2}\) Bogen.}
\end{tabular*}\par\medskip

\noindent\hspace*{\parindent}Im ersten Abschnitt ist zu handeln von folgenden Gegenständen: Nächster
Zweck der Messungen; Eigenthümlichkeit des Terrains; Zielpunkte; Standpunkte;
Instrumente; Beobachtungsmanier; Basis; provisorische Berechnung der Punkte.
\par\smallskip
\noindent\hspace*{.031\linewidth}{\footnotesize IX.}\hfill{\small 51}\hspace*{.128\linewidth}\par

% Printed 402
\clearpage
\noindent\makebox[0pt][l]{\fontsize{13}{16}\selectfont 402}\hfill{\small\textbf{NACHLASS.}}\hfill\null\par
\vspace{2.5\baselineskip}
\begin{center}
{\large\textls[120]{Die trigonometrischen Messungen im Königreich Hannover.}}\par\medskip
\rule{.12\linewidth}{.4pt}\par\medskip
{\large\textls[120]{Einleitung.}}
\end{center}

\par\noindent\hspace*{\parindent}Die erste Veranlassung zu den von mir in dem Königreiche Hannover
in den Jahren 1821---1825 ausgeführten trigonometrischen Messungen war
durch die von dem Herrn Etatsrath \textsc{Schumacher} in den dänischen Staaten
unternommene Gradmessung gegeben. Wenn diese in der Nordspitze von
Jütland ihren natürlichen Endpunkt finden musste, so war sie im Süden einer
Erweiterung fähig, die nur erst am mittelländischen Meere ihre Begrenzung
findet. Zunächst musste eine solche Erweiterung in einer Strecke von mehr
als zwei Breitengraden durch das Königreich Hannover gehen, an dessen südlicher
Grenze die Göttingische Sternwarte die Gelegenheit zu den feinsten
und bequemsten astronomischen Bestimmungen darbot. Die wissenschaftliche
Wichtigkeit einer solchen Unternehmung und die feste Grundlage, welche
dadurch die Geographie des Königreichs und künftige umfassendere trigonometrische
Messungen erhalten mussten, konnten unserm erleuchteten Gouvernement
nicht entgehen, und ich erhielt daher im Jahre 1820 den Auftrag,
diese Fortsetzung der dänischen Gradmessung durch das Königreich Hannover
auszuführen.

\par\noindent\hspace*{\parindent}Der trigonometrische Theil dieser Arbeit wurde im Sommer 1821 am
südlichen Ende angefangen und im Sommer 1823 mit den Messungen auf dem
Michaelisthurm in Hamburg beendigt.

\par\noindent\hspace*{\parindent}Inzwischen war in dem benachbarten Kurfürstenthum Hessen eine trigonometrische
Landesvermessung unter der Leitung des Herrn Professor \textsc{Gerling}
angefangen, bei welcher die Hauptdreiecke mit ausgezeichnet guten Hülfsmitteln
und mit aller erreichbaren Genauigkeit gemessen werden sollten. Eine Verbindung
derselben mit den hannoverschen Dreiecken war daher um so wichtiger,
weil dadurch diese mit den bayerischen Dreiecken in Zusammenhang
kommen mussten, und die in verschiedenen Theilen von Europa ausgeführten
Dreiecksmessungen durch ihre Verknüpfung zu Einem Ganzen in höherer
\par\clearpage

% Printed 403
\noindent\null\hfill{\small\textbf{ZUR HANNOVERSCHEN TRIANGULATION.}}\hfill\makebox[0pt][r]{\fontsize{13}{16}\selectfont 403}\par\medskip
\noindent wissenschaftlicher Beziehung einen vielfach erhöhten Werth erhalten. Jene
Verbindung der hannoverschen und kurhessischen Messungen wurde noch
im Spätjahr 1823 ausgeführt; letztere sind aber seitdem unvollendet geblieben,
obwohl so viel als zu einer nothdürftigen Verbindung der hannoverschen mit
den bayerischen Dreiecken erforderlich war, nemlich die Messung von wenigstens
zwei Winkeln in allen zu der Verbindung nöthigen Dreiecken, im Jahre
1823 vollendet ist.

\par\noindent\hspace*{\parindent}So wie nun hiedurch eine wenigstens vorläufige Verknüpfung der norddeutschen
und süddeutschen Messungen erreicht war, musste es doppelt wichtig
erscheinen, auch eine Verknüpfung mit den grossen unter sich zusammenhängenden
Messungssystemen im Westen zu bewirken, und mein ursprünglicher
Auftrag erhielt deshalb eine Erweiterung, dass ich auch noch einen Übergang
von meinen Gradmessungsdreiecken zu den \textsc{Krayenhoff}schen Messungen ausführen
sollte. Ich führte deshalb in den Jahren 1824 und 1825 zu diesem
Zwecke ein neues Dreieckssystem von Hamburg bis Jever, wodurch also der
Zusammenhang mit den niederländischen, französischen und englischen Dreiecken
bewirkt ist, so dass also schon jetzt alle grossen durch den cultivirtesten
Theil von ganz Europa sich erstreckenden Messungen in der That vorhanden
sind.

\par\noindent\hspace*{\parindent}Da inzwischen die dänische Gradmessung noch unvollendet geblieben
war, so hielt ich für nothwendig, den astronomischen Theil meines Geschäfts
so einzurichten, dass die hannoversche Gradmessung auch als ein abgeschlossenes
Ganzes für sich bestehen konnte.

\par\noindent\hspace*{\parindent}Die Göttingische Sternwarte, welche selbst ein Hauptdreieckspunkt im
System ist, und von der aus alle Dreiecksseiten ihre Orientirung erhielten,
bildete von selbst den südlichen Endpunkt; allein ein in seiner Art einziger
Umstand kam hinzu, der auch die Wahl des nördlichen Endpunkts nicht
zweifelhaft lassen konnte: die inzwischen in Altona errichtete Sternwarte des
Herrn Professor \textsc{Schumacher} liegt nemlich fast genau im Meridian der Göttingischen.
Die von mir im Jahr 1827 an beiden Plätzen mit dem \textsc{Ramsden}schen
Zenithsector gemachten Beobachtungen an 43 Sternen und die daraus
für den Breitenunterschied erhaltenen und sonstigen Resultate habe ich bereits
in einem 1828 erschienenen Werke bekannt gemacht.

\par\noindent\hspace*{\parindent}Dem trigonometrischen Theil meiner Arbeit ist gegenwärtiges Werk ge-
\par\smallskip
\noindent\hfill{\small 51\textsuperscript{*}}\hspace*{.122\linewidth}\par\clearpage

% Printed 404
\noindent\makebox[0pt][l]{\fontsize{13}{16}\selectfont 404}\hfill{\small\textbf{NACHLASS.}}\hfill\null\par\medskip
\noindent widmet. Ich habe dabei aus einem doppelten Grunde für nöthig gehalten,
der Darstellung alle erforderliche Ausführlichkeit zu geben.

\par\noindent\hspace*{\parindent}Bei einer \textls[120]{isolirten} Breitengradmessung von mässiger Ausdehnung,
\textls[120]{die nichts weiter als solche ist und sein soll}, kann man den trigonometrischen
Theil gewissermaassen als untergeordnet betrachten, insofern die Genauigkeit,
deren die astronomischen Beobachtungen fähig sind, doch lange
nicht der bei dem trigonometrischen Theil erreichbaren Genauigkeit entspricht,
und es also nicht so unerlässlich nothwendig ist, in Beziehung auf letztere das
Höchste zu erreichen.

\par\noindent\hspace*{\parindent}Bei der trigonometrischen Vermessung eines Landes ist es dagegen in
mehrern Rücksichten allerdings rathsam, die Genauigkeit in der Bestimmung
der gegenseitigen Lage der Hauptpunkte so weit zu treiben, wie es der Zustand
der Kunst und die Umstände nur zulassen, zumal da es dann in unzähligen
Fällen möglich wird, hinreichend genau abgeleitete Bestimmungen
secundärer Punkte mit äusserst geringer Arbeit und durch Methoden zu gewinnen,
die ohne jene Voraussetzung ins Wilde führen würden. Wenn eine solche
trigonometrische Vermessung isolirt steht, hat freilich ausführlichere Bekanntmachung
ihrer Bestandtheile wenigstens kein allgemeines Interesse. Allein
je mehr die in verschiedenen Theilen von Europa ausgeführten trigonometrischen
Messungen mit einander in Verbindung kommen und nach und nach
sich einem grossen Ganzen nähern werden, desto mehr erhalten die einzelnen
Bestandtheile den Charakter eines kostbaren Gemeinguts von einem für alle
Zeiten bleibenden Werthe, und desto wichtiger wird es, alle wesentlichen Momente
derselben in solcher Vollständigkeit aufzubewahren, dass ihre Zuverlässigkeit
im Ganzen wie im Einzelnen stets geprüft werden könne.

\par\noindent\hspace*{\parindent}Ein zweiter Beweggrund zur Ausführlichkeit lag in der Eigenthümlichkeit
der, sowohl bei den Messungen selbst, als bei ihrer Verarbeitung zu
Resultaten, von mir angewandten Methoden, welche von den sonst üblichen
zum Theil gänzlich verschieden sind, und deren Darstellung ein Hauptzweck
dieses Werks sein sollte. Ohne Zweifel ist die Verbindung einer Darstellung
dieser Methode im allgemeinen, mit einer fortlaufenden Anwendung auf ein
ausgedehntes Messungssystem, das geeignetste Mittel, die Natur derselben in
ihr wahres Licht zu setzen, und denjenigen, welche sich derselben in Zukunft
zu ähnlichen Messungen bedienen wollen, diese Anwendung zu erleichtern.
\par\clearpage

% Printed 405
\noindent\null\hfill{\small\textbf{ZUR HANNOVERSCHEN TRIANGULATION.}}\hfill\makebox[0pt][r]{\fontsize{13}{16}\selectfont 405}\par\medskip
\begin{center}
\textls[120]{Erster Abschnitt.}\par\smallskip
Anordnung der Messungen im allgemeinen.\par\medskip
1.
\end{center}

\par\noindent\hspace*{\parindent}Der Landstrich von Göttingen bis Hamburg ist in seinem südlichen und
nördlichen Theile von sehr ungleicher Beschaffenheit. Jener ist gebirgig, und
die Berge sind auf ihren Gipfeln meistens mehr oder weniger bewaldet, und
die meisten Ortschaften liegen so, dass ihre Thürme eine weite Aussicht entweder
gar nicht oder höchstens nach Einer Seite darbieten. Der nördliche
Theil hingegen ist flach, vielfach mit Waldung durchschnitten, welche die
Benutzung einzelner Anhöhen von geringer Höhe sehr erschwert, und oft
ganz unthunlich macht; und an Ortschaften mit Thürmen, die sich zu Dreieckspunkten
eigneten, fehlt es auf dem grössten Theile dieser Strecke gänzlich.

\begin{center}
2.
\end{center}

\par\noindent\hspace*{\parindent}Unter diesen Umständen liess sich voraussehen, dass auf natürliche Dreieckspunkte
fast gar nicht zu rechnen, sondern an den meisten Dreieckspunkten
entweder eigene Signalthürme zu erbauen, oder auf andere künstliche Mittel
zu ihrer Sichtbarmachung Bedacht zu nehmen sein würde.

\par\noindent\hspace*{\parindent}Ein Hauptumstand in dieser Beziehung ist die Grösse, welche man den
einzelnen Dreiecken zu geben beabsichtigt. Es ist klar, dass bei einer sehr
ins Grosse gehenden, z.~B. einen ganzen Welttheil umfassenden Messung, es,
allgemein zu reden, für die Genauigkeit des Ganzen am vortheilhaftesten sein
würde, die Dreiecke so gross wie nur möglich zu machen, und dasselbe gilt
dann auch für eine Messung von kleinerm Umfang, insofern man sie als einen
Bestandtheil eines solchen grössern Systems betrachtet. Allein diese Behauptung
bleibt nur insofern wahr, als man voraussetzt, die Winkel in den
grössten Dreiecken seien mit derselben, wenigstens mit einer nicht erheblich
geringern Schärfe zu messen, wie die in kleinen Dreiecken, und diese Bedingung
findet freilich bei der Anwendung von Kirchthürmen oder künstlichen
Signalthürmen keinesweges statt, und die Augenblicke, wo dergleichen Gegenstände
in sehr grossen Entfernungen die grösste Schärfe und Sicherheit in den
Messungen verstatten, sind äusserst selten.

\begin{center}
\rule{.24\linewidth}{.4pt}
\end{center}
\clearpage

% Printed 406
\noindent\makebox[0pt][l]{\fontsize{13}{16}\selectfont 406}\hfill{\small\textbf{NACHLASS.}}\hfill\null\par
\vspace{2\baselineskip}
\begin{center}
[2.]\par\smallskip
{\large [Auszüge aus Berichten über die Triangulirung an das hannoversche\\
Cabinets-Ministerium.]}
\end{center}
\bigskip
\begin{center}
[Aus einem Bericht vom 7.~Januar 1822 »über die Arbeiten im Jahre 1821«.]
\end{center}

\par\noindent\hspace*{\parindent}Von den vielfachen Operationen, welche zu einer Gradmessung gehören,
ist die Bildung des Dreiecksnetzes und die Messung der Winkel diejenige,
welche bei weitem die meiste Zeit und Arbeit erfordert. Bei der hannoverschen
Gradmessung muss dies Dreiecksnetz im Norden bei Hamburg sich an
die dänischen Messungen anschliessen: im Süden ist zwar die Göttinger Sternwarte
der eigentliche natürliche Endpunkt der Gradmessung an sich; allein
damit diese auch der weitern Ausdehnung nach Süden fähig werde, ist es zugleich
sehr wesentlich, sie an diejenigen fremden Messungen anzuschliessen,
welche das Königreich Hannover auf der Südseite berühren. In diesem Falle
befinden sich die von der königl.\ preussischen Regierung veranstalteten und
mit grosser Sorgfalt ausgeführten Messungen, so wie gegenwärtig auch in Kurhessen
eine grosse mit aller erreichbaren Genauigkeit auszuführende Triangulirung
beabsichtigt wird.

\par\noindent\hspace*{\parindent}Die erwähnten preussischen Messungen, so weit sie bisher gediehen waren,
wurden mir im vorigen Winter mitgetheilt; ausserdem hatte ich Gelegenheit,
einen Theil der von dem französischen Obersten \textsc{Epailly} im Jahr 1804 u.\ f.
im Hannoverschen, und namentlich im südlichen Theil, gemachten Messungen
zu erhalten. Der Besitz dieser und einiger anderer Hülfsmittel, welche bei
der ersten Auswahl der Dreieckspunkte einige Erleichterungen geben konnten,
sowie die Erwägung, dass manchen kleinen von dem Anfang solcher Operationen
unzertrennlichen Schwierigkeiten und Verlegenheiten immer schneller
und leichter in der Nähe von Göttingen würde abgeholfen werden können,
bestimmten mich, die Triangulirung \textls[120]{auf der Südseite} anzufangen.

\par\noindent\hspace*{\parindent}Schon im Jahre 1820 hatte ich angemessene Einleitungen getroffen, um
mir die erforderlichen Instrumente zu verschaffen\textperiodcentered{} hier erwähne ich nur der-
\par\clearpage

% Printed 407
\noindent\null\hfill{\small\textbf{ZUR HANNOVERSCHEN TRIANGULATION.}}\hfill\makebox[0pt][r]{\fontsize{13}{16}\selectfont 407}\par\medskip
\noindent jenigen, welche sich unmittelbar auf den geodätischen Theil der Gradmessung
beziehen. Zu den eigentlichen Winkelmessungen hatte ich bei \textsc{Reichenbach}
(dessen Werkstatt sein ehemaliger Werkmeister \textsc{Ertel} gegenwärtig ganz übernommen
hat) einen zwölfzölligen Theodolithen bestellt, dessen Vollendung und
Ablieferung auf das Frühjahr 1821 zugesagt war. Einen kleinern Theodolithen
von dem englischen Künstler \textsc{Troughton} hatte ich durch die gefällige Besorgung
des Professors \textsc{Schumacher} bereits in Händen.

\par\noindent\hspace*{\parindent}Eine besondere Vorsorge erforderten die Hülfsmittel, die Dreieckspunkte
in sehr grossen Entfernungen sichtbar zu machen. Da es meine Absicht und
von grösster Wichtigkeit war, die Dreiecke so gross wie möglich zu wählen,
so blieb bei der Beschaffenheit des Landstriches, durch welches sie zu führen
sind, keine Hoffnung, dass viele Kirchthürme als Dreieckspunkte würden benutzt
werden können. Besonders gebaute Signalthürme sind bisher das in
solchen Fällen am meisten angewandte Mittel gewesen: indessen kommen in
der Ausübung nicht selten Fälle vor, wo auch dieses Mittel unzureichend
wird, indem solche Signalthürme (ebenso wie die Kirchthürme) in grossen
Entfernungen, da, wo sie nicht gegen den Himmel, und besonders da, wo sie
sich gegen nahen dunkelfarbigen Hintergrund projiciren, immer sehr schwer
zu sehen, und noch viel schwerer zu beobachten sind\textsuperscript{*}). Andere Beobachter
haben aus diesen und andern Gründen häufig (einige ausschliesslich) die
Winkelbeobachtungen bei Nacht angestellt, indem sie die entfernten Dreieckspunkte
durch grosse \textsc{Argand}sche Lampen mit sehr genau parabolischen Reverberes
sichtbar machen liessen. Freilich haben diese nächtlichen Beobachtungen
wieder andere grosse Schwierigkeiten und Inconvenienzen, und besonders bei
sehr grossen Dreiecken muss gewöhnlich eine gelungene Beobachtung erst mit
vielen vergeblichen Versuchen gleichsam erkauft \mbox{werden.\enspace\mbox{.\ .\ .\ .}}\quad Wenn ich
daher gleich nicht geneigt war, mich dieser Beobachtungsart ausschliesslich
zu bedienen, zumal da meine physischen Kräfte den Beschwerden eines beständigen
nächtlichen Aufenthalts auf meistens hohen und schwer zugänglichen
Bergen schwerlich gewachsen gewesen sein würden, so musste ich mich doch,

\par\smallskip\noindent\rule{.24\linewidth}{.4pt}\par\smallskip
{\footnotesize
\noindent\hspace*{\parindent}\textsuperscript{*}) Meine eigene Erfahrung im vorigen Sommer hat dies vielfach bestätigt.
So habe ich z.~B. während meines ganzen mehr als vierwöchentlichen Aufenthalts auf dem Brocken den auf dem Hils
erbauten \(7\tfrac{1}{2}\) Meilen entfernten Signalthurm nur ein- oder zweimal auf wenige Minuten, die Kirchthürme des
12 Meilen entfernten Hannover auch nicht ein einziges Mal sehen können, ungeachtet die Richtung genau bekannt war.\par}
\clearpage

% Printed 408
\noindent\makebox[0pt][l]{\fontsize{13}{16}\selectfont 408}\hfill{\small\textbf{NACHLASS.}}\hfill\null\par\medskip
\noindent da kein anderes Mittel bisher bekannt war, im voraus gefasst halten, dasselbe
wenigstens in manchen einzelnen Fällen zu gebrauchen, und ich hatte
daher vorläufig drei solcher Lampen bei \textsc{Repsold} in Hamburg und bei \textsc{Körner}
in Jena bestellt. Diese Lampen erhielt ich im Mai 1821, und ihre Wirkung
bei den in schicklichen Entfernungen damit verschiedentlich angestellten Versuchen
hat auch meiner Erwartung entsprochen.

\par\noindent\hspace*{\parindent}Indem mir alle die erwähnten grossen Schwierigkeiten bei Bildung grosser
Dreiecke nach fremden Erfahrungen, noch ehe ich eigene gemacht hatte, vorschwebten,
war ich auf ein ganz neues Mittel bedacht, ihnen abzuhelfen.
Theoretische Untersuchungen hatten mich überzeugt, dass reflectirtes Sonnenlicht
von nur ganz kleinen Planspiegeln hinreichende Kraft habe, um in den
grössten Entfernungen sichtbar zu sein, und sich viel leichter und besser beobachten
zu lassen, als alle Thürme und Signale, ja selbst besser, als mehrere
zusammengestellte \textsc{Argand}sche Lampen bei Nacht. Um diese Idee brauchbar
zu machen, bedurfte es eines besondern Apparats oder Instruments, wodurch
man das reflectirte Sonnenlicht mit grösster Genauigkeit und Sicherheit ununterbrochen
nach jedem beliebigen noch so weit entfernten Punkte lenken
kann. Obgleich ich die Einrichtung zu diesem Zweck im wesentlichen schon
vollständig entworfen hatte, war es doch nicht leicht, dasselbe ausgeführt zu
erhalten, zumal wenn dies durch einen auswärtigen Künstler hätte geschehen
sollen, wo der Vortheil fortwährender mündlicher Berathung bei einzelnen
technischen Schwierigkeiten weggefallen wäre, und die Vollendung daher zum
wenigsten sehr in die Länge gezogen sein würde. Diese Verlegenheit näherte
sich jedoch ihrer Erledigung, als unser geschickter Inspector \textsc{Rumpf}, welcher
während eines grossen Theils des Winters von hier abwesend gewesen war,
gegen Ostern nach Göttingen zurückkehrte, und bald nachher die Arbeit eines
solchen Instruments \mbox{übernahm.\enspace\mbox{.\ .\ .\ .\ .\ .}}

\par\noindent\hspace*{\parindent}Bei der hannoverschen Gradmessung ist es ein besonders wesentlicher und
wichtiger Umstand, dass die ersten von der hiesigen Sternwarte auslaufenden
Dreiecksseiten durch die festen Meridianinstrumente der Sternwarte selbst
auf das genaueste orientirt werden können. Die Aussicht der Sternwarte in
der Richtung des Meridians war zwar ursprünglich weder auf der Nordseite
noch auf der Südseite offen, und wenn diesem Mangel nicht abzuhelfen gewesen
wäre, so wäre nicht allein jener höchst wichtige Vortheil gar nicht
\par\clearpage

% Printed 409
\noindent\null\hfill{\small\textbf{ZUR HANNOVERSCHEN TRIANGULATION.}}\hfill\makebox[0pt][r]{\fontsize{13}{16}\selectfont 409}\par\medskip
\noindent vorhanden gewesen, sondern es wäre dies auch auf immer ein wesentlicher
Radicalfehler bei der Wahl des Platzes der Sternwarte geblieben. Glücklicherweise
war aber die eine Hälfte dieser Schwierigkeit bereits überwunden;
die vorher durch die Gärten vor Göttingen versperrt gewesene Aussicht nach
Norden hatte ich schon im Herbst 1820 geöffnet, und auf einem Berge unweit
Weende ein provisorisches Meridianzeichen errichten lassen. Viel grösser
waren hingegen die Schwierigkeiten auf der Südseite des Meridians, wo eine
dichte hohe drei Stunden entfernte Waldung die Aussicht begrenzte. Dies
Hinderniss musste wo möglich überwunden \mbox{werden.\enspace\mbox{.\ .\ .\ .\ .\ .}}

\par\noindent\hspace*{\parindent}Während dieser Zeit hatte der Inspector \textsc{Rumpf} schon fleissig an dem
oben erwähnten Instrument gearbeitet, welches ich fortan mit dem ihm beigelegten
Namen Heliotrop bezeichnen werde. Allein noch vor dessen Vollendung
war ich auf die Idee gekommen, einen blossen Spiegelsextanten zu
einer Art Viceheliotrop einzurichten, freilich viel unvollkommener, als jenes
Instrument selbst, aber doch bei geschickter Behandlung gleichfalls brauchbar.
Die damit schon auf Entfernungen von beinahe 2 Meilen angestellten Versuche
bestätigten die enorme Kraft des reflectirten Sonnenlichts fast über
meine \mbox{Erwartung.\enspace\mbox{.\ .\ .\ .\ .\ .}}
\par\bigskip

\begin{center}
[Aus einem Bericht vom 31.~Januar 1823 »über die Arbeiten der Gradmessung im Jahre 1822«.]
\end{center}

\par\noindent\hspace*{\parindent}Im Jahre 1821 war die Triangulirung, als der ausgedehnteste Theil des
ganzen Geschäfts, eingeleitet, und die Messungen an fünf Dreieckspunkten
vollführt, nemlich in der Göttinger Sternwarte, beim Meridianzeichen, auf dem
Hohehagen, Hils und Brocken. Behuf der weitern Fortsetzung waren ferner
vier weiter nördlich liegende Punkte ausersehen, nemlich Lichtenberg, der
Deister, der Wohlenberg im Amte Gifhorn, und ein Berg bei Brelingen
in der Amtsvoigtei Bissendorf. Ich selbst hatte jedoch diese Punkte noch
nicht besucht, und es blieb bei einigen derselben noch problematisch, ob sie
sich zu Dreieckspunkten qualificiren würden: dies musste erst entschieden
werden, ehe der Plan zu den ersten Arbeiten des Jahrs 1822 gemacht werden
konnte.

\par\smallskip
\noindent\hspace*{.030\linewidth}{\footnotesize IX.}\hfill{\small 52}\hspace*{.137\linewidth}\par
\clearpage

% Printed 410
\noindent\makebox[0pt][l]{\fontsize{13}{16}\selectfont 410}\hfill{\small\textbf{NACHLASS.}}\hfill\null\par\medskip
\par\noindent\hspace*{\parindent}Allein dies war nur der kleinste Theil der nothwendigen Präliminar-Untersuchungen.
Die Operationen näherten sich nun der Lüneburger Heide,
einem ganz flachen Lande, wo der Mangel dominirender Punkte und die fast
unzählbaren grössern und kleinern Holzungen, welche es schachbrettartig bedecken,
die Bildung von einigermaassen beträchtlichen Dreiecken ausserordentlich
erschweren. Ich kannte diese Schwierigkeiten bereits aus den Berichten
des französischen Obersten \textsc{Epailly}, der im Jahr 1804 und 1805 die französischen
Messungen im Kurfürstenthum Hannover geleitet hatte: dieser Ingenieur
hatte die Schwierigkeiten des Terrains für so gross angesehen, dass er
die Bildung eines Dreieckssystems von der Aller bis zur Elbe für unmöglich
erklärt, und daher die Verbindung auf einem ungeheuer grossen Umwege,
nemlich durch Dreiecke längs der Weser bis zu ihrer Mündung und dann
wieder die Elbe herauf bis Hamburg, effectuirt hatte, ein Verfahren, das
höchstens als Nothbehelf bei einer Landesvermessung, aber durchaus nicht
bei einer Gradmessung zulässig sein könnte.

\par\noindent\hspace*{\parindent}Diese Umstände machten vor dem Anfange der eigentlichen Messungsoperationen
eine Recognoscirungsreise nothwendig. Diese Reise, bei welcher
ich von meinen drei Gehülfen nur den Hauptmann \textsc{Müller} mit zuzog,
beschäftigte mich vom 28.~April bis 1.~Junius, und ich führe, das Detail der
mühsamen Untersuchungen hier übergehend, nur die Hauptresultate derselben
an.

\par\noindent\hspace*{\parindent}Der Brelingerberg wurde zur Verbindung unbrauchbar befunden; dies
wurde aber mehr als ersetzt durch die glückliche Entdeckung, dass zwei neue
noch nördlicher liegende Punkte, ein hoher Acker bei Garssen und der Falkenberg,
jener fast eine Meile nordöstlich von Celle, dieser eine Meile nordwestlich
von Bergen entfernt, sich beide unmittelbar mit Lichtenberg und mit dem
Deister verbinden liessen: dadurch wurde der Wohlenberg überflüssig, und die
ersten Messungsarbeiten waren nun bestimmt und sicher festgesetzt. In Rücksicht
auf das weitere Fortschreiten nach Norden fand ich allerdings die
Schwierigkeiten so gross, wie ich erwartet hatte; jedoch war es mir auch gelungen,
gleichsam im Herzen der Heide die Ausführbarkeit zweier guter Dreiecke
zwischen den vier Punkten Falkenberg, Hauselberg (in der Amtsvoigtei
Hermannsburg), Wulfsode (im Amt Ebstorf) und Wilsede (an der äussersten
südwestlichen Grenze des Amts Winsen an der Luhe) festzustellen. Es zeigte
